\documentclass[11pt,a4paper]{article}
\usepackage[utf8]{inputenc}
\usepackage[T1]{fontenc}
\usepackage{lmodern}
\usepackage{microtype}
\usepackage[margin=1in]{geometry}
\usepackage{amsmath,amssymb,bm}
\usepackage{booktabs}
\usepackage{hyperref}
\usepackage{parskip}
\title{Weekly Meeting Update:\\ HPO Campaign, Benchmarking, and Gradient Routing (GR-Share)}
\author{Niklas Nertinger}
\date{May 8, 2026}
\begin{document}
\maketitle
\section*{Overview}
This document outlines the current status of the VarShare multi-task reinforcement learning project. We recap the mathematical formulation of our proposed \textit{Gradient Routing} method (working title: \textbf{GR-Share}) and describe the current large-scale cluster HPO campaign. We report empirical computational profiles, establish our design alignment with the MTBench benchmark, and outline our multi-task diagnostic metrics. Finally, we discuss our hyperparameter optimization (HPO) pruning strategy and our ongoing work extending GR-Share to the Soft Actor-Critic (SAC) framework.
\section{Introduction to Gradient Routing (GR-Share)}
The central challenge in Multi-Task Reinforcement Learning (MTRL) is balancing joint feature acquisition (\textit{constructive interference}) with task-specific conflict resolution (\textit{destructive interference} or \textit{gradient conflict}). 
We introduce task-specific additive parameter adapters $\mu_t$ alongside a shared core backbone $\theta$, defining the effective parameters for task $t$ as:
\begin{equation}
    w_t = \theta + \mu_t
\end{equation}
Let $\ell_t(w_t)$ denote the task-$t$ loss. The raw task gradient is:
\begin{equation}
    g_t := \nabla_{w_t} \ell_t(w_t)
\end{equation}
To protect the shared backbone from destructive interference, we apply gradient surgery (e.g., PCGrad projection) to the raw task gradients $\{g_t\}_{t=1}^T$ to obtain conflict-resolved modified gradients $\tilde{g}_t$. The backbone update is formed by averaging:
\begin{equation}
    \bar{g}_{\theta} := \frac{1}{T}\sum_{t=1}^T \tilde{g}_t
\end{equation}
Rather than discarding the gradient components removed during surgery, GR-Share \textit{routes} this task-specific residual information back into the adapters. We define the routed residual as:
\begin{equation}
    r_t := g_t - \bar{g}_{\theta}
\end{equation}
The task-specific adapters are then updated using this residual scaled by routing strength $\alpha \ge 0$ and pulled back toward zero via explicit shrinkage coefficient $\lambda > 0$:
\begin{equation}
    \bar{g}_{\mu_t} := \alpha r_t + \lambda \mu_t
\end{equation}
Using standard gradient descent, the resulting updates are:
\begin{align}
    \theta^{+} &= \theta - \eta \bar{g}_{\theta} \\
    \mu_t^{+}  &= (1 - \eta\lambda)\mu_t - \eta\alpha(g_t - \bar{g}_{\theta})
\end{align}
Under this update rule, the shared backbone absorbs joint generalist knowledge, while task-specific discrepancies that conflict with the consensus are safely diverted into explicitly regularized, task-specific pathways.
\section{Cluster Campaign & Experimental Configuration}
We are currently executing a comprehensive, fully parallel hyperparameter optimization (HPO) and benchmarking campaign on the cluster. The study consists of \textbf{7 distinct baseline methods} plus the Oracle upper bound (8 total baseline runs) along with GR-Share:
\begin{itemize}
    \item \textbf{Baselines:} \texttt{Shared}, \texttt{Independent}, \texttt{PCGrad}, \texttt{CAGrad}, \texttt{PaCo}, \texttt{Soft Modularization}, and \texttt{MOORE}.
    \item \textbf{Standardized Capacity Constraints:} To ensure a scientifically rigorous and fair comparison, all modular/multi-expert baselines are constrained to match our standard 3-layer, 400-width backbone. For any method utilizing $K$ experts or modules, we dynamically scale its hidden dimension according to:
    \begin{equation}
        D = \text{int}\left(\frac{400}{\sqrt{K}}\right)
    \end{equation}
    This limits all models to approximately equal active parameter counts ($\sim$330k parameters), preventing larger baselines from succeeding purely due to capacity advantages.
\end{itemize}
\section{Computational Speed vs. Sample Efficiency}
Empirical profiles from our cluster timing runs at the 100,000 environment step mark reveal the following computational footprint for GR-Share (Routing) relative to the baselines:
\begin{itemize}
    \item \textbf{Vs. Lightweight Baselines:} GR-Share (637.15s) is \textbf{$\sim$5.3 times slower} than basic architectures like \texttt{Shared} (111.64s), \texttt{Independent} (122.19s), and \texttt{MOORE} (129.68s).
    \item \textbf{Vs. Gradient Surgery Baselines:} GR-Share is \textbf{$\sim$1.65 to 1.8 times slower} than standard PCGrad (386.95s) and CAGrad (356.79s).
    \item \textbf{Vs. Base VarShare:} GR-Share is \textbf{$\sim$2.7 times slower} than the base deterministic formulation (\texttt{det\_base} at 238.03s).
\end{itemize}
\textit{Note: These ratios are achieved after bringing down the compute time of routing by a factor of $\sim$2.5 through aggressive code optimizations, including vectorizing pairwise projection operations, caching flattened tensors, and streamlining the gradient routing loops.}
Importantly, our goal with GR-Share is not wall-clock execution speed, but \textbf{sample efficiency}. Preliminary learning curves are highly encouraging, indicating that the sample-efficiency gains from protecting the shared backbone substantially outweigh the computational overhead.
\section{The MTBench Reference & PPO Rationale}
While the vast majority of modern MTRL literature (such as Soft Modularization or PaCo) defaults to Soft Actor-Critic (SAC) as its underlying algorithm, we built our architecture natively on Proximal Policy Optimization (PPO). 
To provide a peer-reviewed, defensible argument for this choice, we align our configuration with the recent \textbf{MTBench} paper. MTBench establishes a rigorous benchmarking suite using PPO on MetaWorld. Although direct numerical comparisons are limited due to their focus on massive parallel environments and wall-clock optimization, their repository provided a valuable resource for:
\begin{enumerate}
    \item Properly implementing and stabilizing multi-task PPO updates.
    \item Selecting a defensible, standardized architecture (3 layers of 400 units).
    \item Sourcing robust default hyperparameters, which we can reference in our work.
\end{enumerate}
\section{Metric Tracking \& The Generalist-Specialist Validation}

To deeply probe the mechanics of joint training and capacity allocation, we track a suite of non-trivial diagnostic and structural metrics during our campaigns:
\begin{itemize}
    \item \textbf{Critic Temporal Difference (TD) Error Variance (\texttt{diagnostics/td\_error\_variance}):} Serves as a primary diagnostic for multi-task critic failure. When a shared critic is confused by competing reward scales across tasks, TD error variance spikes. Tracking this helps us identify if our architecture successfully stabilizes the critic.
    \item \textbf{Critic Explained Variance (\texttt{diagnostics/explained\_variance}):} Measures the stability of the critic's value predictions relative to actual task returns.
    \item \textbf{Gradient Norm (\texttt{diagnostics/grad\_norm}):} Tracks the pre-clipped $L_2$ norm of the gradients to directly quantify the extent of active gradient conflicts within the shared layers.
    \item \textbf{Global Sharing Ratio:} Formally defined as:
    \begin{equation}
        \text{Sharing Ratio} := \frac{\|\theta\|_2}{\|\theta\|_2 + \frac{1}{T}\sum_{t=1}^T \|\mu_t\|_2}
    \end{equation}
    This metric provides a direct, continuous reading of the model's reliance on generalist versus task-specific pathways.
    \item \textbf{High-Resolution Structural Mapping:} We track the $L_2$ norms and sharing ratios of task-specific adapters decomposed both \textit{per-task} and \textit{per-layer} simultaneously. This reveals whether specialization is localized in early state/visual layers or deeper policy layers.
\end{itemize}

Our core validation of the GR-Share architecture relies on a custom evaluation ablation: \textbf{The Generalist vs. Specialist Ablation}:
\begin{itemize}
    \item \textbf{Mechanism:} During evaluation rollouts, we artificially set all task-specific adapters to zero ($\mu_t = 0$) and evaluate the agent using \textit{strictly} the shared backbone weights $\theta$.
    \item \textbf{Objective:} This ablation mathematically tests whether the shared backbone has successfully isolated and learned a generalist baseline policy (e.g., navigating toward target objects) while leaving the adapters to handle fine-grained, task-specific details (e.g., grasping and inserting). Successful demonstration of this behavior provides irrefutable proof of the intended weight-space separation.
\end{itemize}

\textbf{Open Question for the Team:} \textit{Beyond traditional metrics like task success or parameter norms, does anyone have experience tracking or designing specific metrics that effectively capture or reveal multi-task coordination bottlenecks, representation drift, or parameter alignment in deep RL?}
\section{HPO Hyperparameter Search \& Pruning Strategy}
To efficiently explore the joint hyperparameter space of the shared base and method-specific modules, we employ an automated Hyperparameter Optimization (HPO) campaign with active pruning:
\begin{itemize}
    \item \textbf{Pruning Methodology:} We utilize the Asynchronous Successive Halving Algorithm (ASHA) to dynamically terminate underperforming trials. Pruning decisions are based on intermediate rolling evaluation rewards rather than raw instantaneous rewards, protecting slow-starting but high-potential trials from premature termination.
    \item \textbf{Exploration Efficiency:} This pruning framework allows us to allocate our cluster budget dynamically, focusing computational resources on the most promising parameter regions and significantly scaling the number of configurations we can explore.
    \item \textbf{Open Question for the Team:} \textit{Given that RL training curves are notoriously non-monotonous and highly unpredictable during early exploration, has anyone had success tuning or overriding pruning thresholds in HPOs for RL without killing slow-starting but ultimately superior runs?}
\end{itemize}
\section{Next Step: Extending GR-Share to SAC}
We are actively working on implementing and optimizing GR-Share natively within Soft Actor-Critic (SAC). Delivering robust empirical results across both PPO (on discrete and continuous setups) and SAC (on continuous setups) in the final paper will provide an exceptionally strong, framework-agnostic argument for the general utility of the Gradient Routing paradigm.
\end{document}